\section{Proof of Theorem~\ref{thm:goldbach}}\label{sec:proof-theorem-1}

\subsection{One signed assembly}

Put $L=\log N$, $\rho=1+\tau$ and
$W_{\rm norm}=4(1-\varepsilon_A)U_C$.
This scalar differs from the pointwise weight.
The signed source is
\[
\begin{split}
W_{\rm norm}D_{\rm src}(N)={}&3S_1+S_2-2S_3-S_4-S_5-S_6-S_7\\
&+S_8+S_9+S_{10}-2S_{11}-S_{12}-S_{13}-S_{\rm sw}.
\end{split}\tag{6.1}
\]

\begin{theorem}[Signed assembly]\label{thm:signed-assembly}
Fix one actual construction with
$u=1/500$, $\alpha=100/1149$, $\beta=50/309$,
$0<\sigma\le10^{-10}$, $\tau=\sigma^9$ and
$\ell_0=53/412-1/20000>0$.
Assume four conditions.
\begin{enumerate}
\item The nineteen row chains and outer chains cover (6.1).
They retain exact source maps, multiplicities, half-open domains,
selectors and restrictions.
\item The source, AA37, comparator, detection, boundary, Class-I,
Class-II, measure and profile estimates below hold.
They retain actual domains, signed modulus sums and norms.
Uniformity holds after the displayed family count.
\item Directed endpoints bound their specified profiles.
S1, S2-T8/4 and G1 enter jointly, once each.
The S1-3 prefactor comparison is separate from its dynamic cutoff.
\item Each fixed loss enters once under the paired tail convention.
The AA37 perturbation fits $g_{\rm AA}$.
The explicit selector, cutoff and peel bounds remain.
Parameters follow the order below.
\end{enumerate}
There are finite signed terms $a_j(N)$ and slacks $s_k(N)\ge0$ with
\[
\begin{gathered}
D_{\rm src}(N)=B_{\rm main}+\sum_j a_j(N)+\sum_k s_k(N),\\
D_{\rm src}(N)\ge B_{\rm main}-F_{\rm pers}-E_{\rm van}(N),
\qquad E_{\rm van}(N)\longrightarrow0.
\end{gathered}\tag{6.2}
\]
Hence the fixed-$\epsilon_A$ threshold holds eventually when
\[
D_{\rm lower}\le B_{\rm main}-F_{\rm pers},\qquad
D_{\rm lower}>T.\tag{6.3}
\]
\end{theorem}

\begin{proof}
The signed row coefficient is included in $Q_{r,k}$.
All quantities below are in D units.
\begin{longtable}{@{}p{0.13\textwidth}p{0.82\textwidth}@{}}
\toprule
Stage & Quantity \\\midrule
$Q_{r,0}$ & Actual signed master source. \\
$Q_{r,1}$ & Declared row source after its exact identity. \\
$Q_{r,2}$ & Low/upper cuts, source tags, squarefree/coprime restoration. \\
$Q_{r,3}$ & Exact selector, dynamic-cutoff and peel partition, with main-comparison error terms. \\
$Q_{r,4}$ & Actual-level AA37 main term, fixed/asymptotic perturbations and remainder. \\
$Q_{r,5}$ & Exact comparator decomposition. \\
$Q_{r,6}$ & $\Lambda/\log$ detection, proper powers and LOG freezing. \\
$Q_{r,7}$ & Boundary, II, I-1 and I-2 partition. \\
$Q_{r,8}$ & Continuum mass after prime-measure/Euler conversion. \\
$Q_{r,9}$ & Sign-correct specified profile after domination. \\
$Q_{r,10}$ & Directed endpoint with the paired tail convention. \\
\bottomrule
\end{longtable}
Finite telescoping yields
\[
Q_{r,0}=Q_{r,10}+\sum_{k=1}^{10}(Q_{r,k-1}-Q_{r,k}).\tag{6.4}
\]
For S2 we use (3.3)--(3.4):
\[
\begin{gathered}
\mathcal F=\Delta\cap\{h\le K\},\quad
\mathcal B=\Delta\cap\{h>K\},\quad R1=\Delta\cap\{g>249/250\},\\
\Delta=PF\ \dot\cup\ (\mathcal F\cap R1)\ \dot\cup\ \mathcal B.
\end{gathered}
\]
Here $h$ is evaluated at inherited labels, and $K$ differs from $K_0$.
S10 Ret/Drop, S5 low/top and S4/S7 A/B retain their half-open endpoints.
The modulus masks are $1-K_T$, $K_T(1-K_P)$ and
$K_TK_P\mathbf1_{\rm rough\ square}$.
They give, before the signed row coefficient,
\[
R_{\rm AA}=S_{\rm slot}+NC_{\rm sf}
-\sum\Xi_{\rm union}\Delta_P
+\sum\Xi_{\rm allowed}(\Delta_0-\Delta_P).\tag{6.4a}
\]
The first two masks split $\Xi_{\rm union}$ disjointly,
and the third gives $\Xi_{\rm allowed}$.
NC uses the remaining squarefree moduli.
A failure of both conditions belongs to the tail.

For $O\in\{\mathrm{S1},\mathrm{S2\text{-}T8/4},\mathrm{G1}\}$,
put $Q_{O,9}=\sum_{r\in O}Q_{r,9}$ and $Q_{O,10}=E_O$.
Then
\[
\begin{split}
\sum_{r\in O}Q_{r,0}={}&Q_{O,10}+(Q_{O,9}-Q_{O,10})\\
&+\sum_{r\in O}\sum_{k=1}^{9}(Q_{r,k-1}-Q_{r,k}),
\qquad Q_{O,9}-Q_{O,10}\ge0.
\end{split}\tag{6.4d}
\]
Signed differences are kept and favourable slacks are discarded.
Each tuple is assigned to its first failed condition.
The low direct primes, the upper G1 cut, the excluded factors and the squarefree/coprime restoration are removed.
The source is then retained.
The unique D1 Lemma~\ref{lem:source-replacement} and the G1
cut-measure comparison prove these changes.

A box meeting a source face is a boundary box.
Otherwise the largest eligible lower-endpoint exponent is chosen for II.
Ties are broken by the fixed variable index, and then I-1 and I-2 are used.
The DIRECT preparation lemmas verify (SH1)--(SH5)
\cite[\S3.11, (3.98)--(3.101)]{Z1}.
Every G1 coordinate exceeds $N^\alpha$,
and its two Class-I families are empty by the product-source proof.
The tables bound every adverse difference,
and the fixed-parameter argument gives $E_{\rm van}\to0$.
This proves (6.2).
\end{proof}

\subsection{Sources, local inputs and directed endpoints}

The map fixes each sign and domain.
Positive terms use lower endpoints, and negative terms use upper endpoints.
\begin{longtable}{@{}p{0.24\textwidth}p{0.71\textwidth}@{}}
\toprule
Contribution & Source and endpoint direction \\\midrule
$3S_1$ & (2.19): coefficients $(3,-3,3,-3)$, orbits $1,2,6$; one joint $\tfrac34\mathcal S_1$ from (A.4). \\
$S_2$ & Four terms; $PF$, $\mathcal F\cap R1$, $\mathcal B$ in (3.4). Whole-$K$ uses II; strict complement uses BV-R1. One lower $d_{T8}$ from (A.7). \\
$-2S_3,S_8,S_9$ & Absolute single-prefix profiles with level reduced by $u=1/500$. \\
$-S_4,-S_7$ & A/B partitions within each row; upper endpoints. \\
$-S_5$ & S5-low on $[1/4,t_*)$ and disjoint BV-top; upper endpoints. \\
$-S_6$ & Separate BV-S6 upper estimate. \\
$S_{10}$ & Ret $\dot\cup$ Drop, with its directed lower estimate. \\
$-2S_{11}-S_{12}-S_{13}$ & G1 chain (4.18), separate cutoff peel; source weights $(2,1,1)$ occur inside $A_G^{\rm cut}$. Child signs $(-,+,-,+)$ give upper/lower/upper/lower profiles; one joint $E_{\rm G1}$. \\
$-S_{\rm sw}$ & O1 (1.9), (1.12), upper $2H_{L,\rm upper}$. \\
\bottomrule
\end{longtable}
The master coefficients are
$(3,1,-2,-1,-1,-1,-1,1,1,1,-2,-1,-1,-1)$.

D1/G1 replace the sources before Buchstab/HB, with their signed corrections.
DIRECT uses Lemmas~\ref{lem:direct-preparation} and~\ref{lem:source-replacement}.
G1 uses (4.14), the monotonicity $N^{1/2}\to N^\mu$ and its cut-measure comparison.
Sections~2--4 verify the actual source and local hypotheses of
\cite[Theorem~3.2, (3.13)]{Z1} and the packing of \cite[Theorem~4.1]{Z1}.
The complete-weight prime-source theorem has a different coefficient family.
For AA37 we apply \cite[Theorem~2.6, (2.30)--(2.33b)]{Z1} with
\[
D=L_c,\quad z=z_k,\quad R=1,\quad
\epsilon_s=\sigma,\quad\delta_s=\tau=\sigma^9.
\]
Section~1 supplies the density and sieve mass.
Sections~2--4 supply the levels and cutoffs.
The three terms of this theorem are $c_kX_kV_N(z_k)M_k(s_k)$,
$-|c_k|X_kV_N(z_k)E_{s,k}$ and $a_kR_{{\rm AA},k}$.
Outer contributions and adverse source differences complete (6.1).

The nineteen applications are indexed by $k=(r,b,\mathcal C)$,
that is, by row, parent-prefix cell and signed remainder.
The remainder $R_{AA,k}$ is taken before the exterior coefficient $a_k$.
It includes the source coefficient and the AA37 sign once,
and it excludes the orbit $M_k$.
For G1, $a_k=(2A_R)\nu_k$, with $\nu_k=+1$ for upper remainders
and $-1$ for lower remainders.
The quantities $C_k,NC_{{\rm sf},k},\Xi_{U,k},\Xi_{A,k}$ are formed before this sign.
Put $\gamma_k=C_k/C_{{\rm slot},k}$.
By \cite[(2.32)--(2.33)]{Z1}, (3.8) and (4.22)--(4.26),
\[
\begin{split}
\frac1{W_{\rm norm}}\sum_k a_kR_{AA,k}
={}&\frac1{W_{\rm norm}}\sum_k\frac{a_kC_{{\rm slot},k}}{M_k}
                  \sum_q\gamma_k(q)\Delta_{P,k}(q)\\
&+\frac1{W_{\rm norm}}\sum_k a_kNC_{{\rm sf},k}
-\frac1{W_{\rm norm}}\sum_{k,q}a_k\Xi_{U,k}(q)\Delta_{P,k}(q)\\
&+\frac1{W_{\rm norm}}\sum_{k,q}a_k\Xi_{A,k}(q)
                    (\Delta_{0,k}-\Delta_{P,k})(q).
\end{split}\tag{6.4b}
\]
This identity is finite at fixed $N$.
Here $r(\mathcal A_k^{\rm pl},d)=\Delta_0(b_kd)$ by (2.17c) and (4.21),
with the AA37 remainder convention of \cite[(2.32)--(2.33)]{Z1}.
For S2 we use the per-block terms underlying (3.8).
The DIRECT comparator differences vanish, while G1 keeps their bounds.
HB is applied to the first term. Its Class-II factor is
\[
\frac{a_kh_{k,B}}{M_kW_{\rm norm}}S^{\rm PSDT}_{k,B}.\tag{6.4c}
\]
Here $h_{k,B}$ comes from the HB/log/box expansion of the same source.
Division by $C_{\rm slot}$ leaves this factor once.
The other terms enter the boundary, proper-power, LOG, I-1 or I-2 classes.
Each block keeps its tested measure, grid, coefficients, cutoff and normalization.

The hypotheses are checked as follows.
\begin{longtable}{@{}p{0.20\textwidth}p{0.75\textwidth}@{}}
\toprule
Input & Actual-source check \\\midrule
C1--C5, LEN & H1--H5 of \cite[Theorem~3.2]{Z1}; Sections~2--4 check coefficients, packing, shift, norms and scales. $a=N$, $N/L^{31}\le X\le N$, $q\le X^{0.6}$, $C_{\rm box}=C_b=334$, $K_{\rm scale}=31$. Direct preparation gives tested measure and integer intervals; G1-LOCAL uses the plain cut measure, fixed majorant and doubled support. \\
H6 & Source-wide freezing and boundary bounds, separately; Section~4.4 includes proper powers. \\
H7 & Both direct families have exponent $2012$; G1 has envelope $K_0=8042$. The finite sum preserves it, including representation multiplicities \cite[(3.91)]{Z1}. \\
Sources and masks & Sections~2.1--2.5, 3.1--3.2, 4.1--4.2; D1/G1 replacement; (6.4a)--(6.4b). \\
Profiles and cutoffs & Sections~3.5, 4.6; (2.22), (2.38)--(2.40), (3.21), (4.45)--(4.57). Squarefree/II: (4.64)--(4.65), (4.68)--(4.69). BV: Sections~3.5, 5.2--5.5. \\
Joint endpoints & Appendix A and the next two comparisons; signs from (6.1). \\
Fixed losses & (1.26)--(1.29), (6.6), parameter order below. \\
\bottomrule
\end{longtable}
The repeated-prime bound is
$232\Lambda_\sigma C_\sigma N^{-\ell_0\sigma^2}(1+L)^4$,
proved in (B.1)--(B.7).
Thus the four assembly hypotheses hold for these families.

For G1 set
\[
\begin{gathered}
E_{\rm G1}=L_{\rm G1}-\delta_{\rm G1},\qquad
\delta_{\rm G1}=4.98\cdot10^{-54},\\
L_{\rm G1}-R_{\rm G1}
=\frac{500103}{\scriptstyle100433627766186892221372630771322662657637687111424552206336},\\
L_{\rm G1}-R_{\rm G1}<4.98\cdot10^{-54},\qquad E_{\rm G1}<R_{\rm G1}.
\end{gathered}
\]
Here $R_{\rm G1}$ sums the four exact row lower endpoints.
We insert $L_{\rm G1}$ from (4.84) in $B_{\rm main}$ and
$\delta_{\rm G1}$ in $F_{\rm pers}$, once each.
Their difference is $E_{\rm G1}=L_{\rm G1}-\delta_{\rm G1}$ in (4.85).
The signed allocation underlying the row endpoints is (A.9).
The row bounds prove (6.4d).
S1 and S2 enter only through their joint endpoints.
Sections~3.1--3.5 assign all four S2 terms, the three pieces, the selector and the tail.

The S1-3 coefficient is $-3$ in (2.19).
Its single $p_3\to E_3$ replacement is (2.20).
The single comparison with $\Phi$ is (2.21)--(2.22).
Thus $E_{1,\rm pref}$ occurs once, at $Q_{\rm S1\text{-}3,9}$, with adverse value $0$.
Its cutoff cost lies inside $62500\tau$, and S2-3 has its own $E_3$.
The correction $E_{1,fF}$ changes S1-1 from $f$ to $F$ instead.

\subsection{The complete error sum}

Every entry below is a nonnegative bound in D units or a proved zero.
The finite logarithmic exponents $C$ are chosen after the earlier constants.
They may differ by row.
\begingroup\small
\setlength{\tabcolsep}{3pt}
\renewcommand{\arraystretch}{1.15}
\begin{longtable}{@{}p{0.17\textwidth}p{0.4\textwidth}p{0.34\textwidth}@{}}
contribution & Scope and exact error term & Bound \\ \hline
\endfirsthead
contribution & Scope and exact error term & Bound \\ \hline
\endhead
Source identities & The 19 signed identities, (6.4b) & $0$ \\
Source restoration & Squarefree and $(n,N)=1$ restoration on 15 direct sources, before AA37 & $O_u(N^{-u/(1+\tau)}L_N^3)$ \\
Direct lower cut & Original $p\le N/L^{30}$ tail, counted before boxes & $O_u(L^{-28})$ \\
Switched upper cut & Actual $e>N(1-L^{1-C_b})$ strip & $O_\sigma(L^{3-C_b}+L^2/N)$ \\
Source endpoint & The original $p<(1-\epsilon_A)N$ source face & $0$; straddling boxes belong to boundary \\
Exceptional source factors & Direct $p\mid N$ or switched tagged factor dividing $N$, at first removal & $O_\sigma(L^3/N+N^{-\alpha/2+o(1)})$ \\
S2 selector & Extra selected cells and dynamic comparison & $275009680\tau$ \\
S10 selector & Dropped whole cells & $7.160\tau$ \\
S1 dynamic cutoff & S1-3 cutoff and dynamic comparisons; prefactor excluded & $62500\tau$ \\
G1 cutoff peel & G1-3 four-prefix peel, never HB-expanded & $145664566.4567\tau$ \\
G1 profile boundary & Matched-profile modulus-label strips; peel separate & $0$ adverse \\
$g_{\rm BV}$ & S5-top $E_{5,\kappa}=4\cdot10^{-8}$ plus an extra allowance $g_{R1}=5.71\cdot10^{-6}$; this allowance is not the corner bound: $d_{T8}$ already includes the R1 corner at $\kappa=10^{-6}$ & $0.00000575$, once \\
$\delta_{\rm G1}$ & Joint-versus-row G1 comparison above & $4.98\cdot10^{-54}$, once \\
$g_\tau$ & Replaces the four selector/cutoff/peel bounds above & $2\cdot10^{-24}$, once \\
Fixed sieve perturbation & Integrated absolute AA37 error on all rows & $g_{\rm AA}=0.00002853328$, once \\
Asymptotic sieve perturbation & Remaining AA37 perturbation & $o_\sigma(1)$; rate unspecified \\
Squareful tail, repeated prefix, allowed rough repetition & Disjoint masks of (6.4a): first forbidden tail equality; first prefix repetition after allowed tail; retained rough repetition. Tests: $\Delta_P$, $\Delta_P$, $\Delta_0-\Delta_P$. & Theorem~\ref{thm:repeated-primes}: $O_\sigma(N^{-\ell_0\sigma^2}L^4)$ separately and together. Third term $0$ on DIRECT. \\
Squarefree normalization & Retained squarefree full moduli only & Direct $0$ on the operative tested-prime source; G1 $O_\sigma(N^{-\alpha}L^5)$ \\
Comparator density & $\omega_b(p^k)=p/(p-1)$ for every $k$; discrepancy in the preceding two rows & $0$ additional \\
Direct prime powers & Tested-side proper prime powers & $O_\sigma(N^{-1/2+o(1)})$ \\
Switched prime powers & At most four tested coordinates, each at least $N^\alpha$ & $O_\sigma(N^{-25/1149+o(1)}L^5)$ \\
Logarithm freezing & Exact $\log^{-1}$-freezing difference over all disjoint boxes & G1 $O_\sigma(L^{166-C_b})$; direct $O_\sigma(L^{20-C_b})$ \\
Logarithmic identity & Freezing is selected; Laplace truncation is unused & $0$ \\
Switched boundary & All four switched rows, one grouped contribution & $3.6\times10^{10}C_\sigma N_{\rm lab}(\sigma)(1+L)^{163}L^{2-C_b/2}$ \\
Direct boundary & All 15 direct rows, one grouped contribution & $3.94\times10^6C_\sigma N_{\rm lab}(\sigma)(1+L)^{16}L^{2-C_b/2}$ \\
Class I-1, Class I-2 & First-failure direct classes; count $K_I$ before choosing $A_i$ & $O_\sigma(L^{K_I+2-A_i})$ \\
Class II & All actual blocks with absolute exterior coefficients & Both terms of (6.5), once \\
Prime-measure and Euler-product conversion & PNT, Li, radial/product prime measure, sieve product and large $p\mid N$ factors & $o_\sigma(1)$ each; \S{}3.5, \S{}4.6 (4.76)--(4.78), direct-BV terms in \S{}5. O1 (5.33)--(5.43) enters $E_{\rm Ssw}$. \\
Profile comparison & Signed domination of the specified profile at the actual level & $0$ adverse by the row comparisons in \S\S{} 2--4; favourable slack discarded \\
$E_{1,\rm pref}$ & S1-3 prefactor, (2.20)--(2.21), at $Q_{\rm S1-3,9}$; dynamic cutoff excluded & $0$ adverse, once \\
$E_{1,fF}$ & S1-1 correction; the S2 correction is already inside its joint endpoint & One error $10^{-390}$ \\
S2 infinite tail & Joint S2 A=1 tail beyond the finite prefix & One error $1/20000$, already included when $d_{T8}$ is used \\
Directed enclosure & Directed numerical enclosure of the specified profile & $0$ additional \\
R1 BV remainder, S5 BV remainder, S6 BV remainder & Three actual BV/sieve bridge errors & $o(1)$, including $O(L^{-1/3})$, $O_A(L^{2-A})$ and the recorded $O(\log L/L)$ measure conversion \\
S6 level increment & Selected S6 $\kappa=10^{-6}$ increment over the nominal S6 baseline & One error $0.000006484294227956901$ \\
Switched-excess error & Full O1 (5.33)--(5.43): source exceptions, outer boxes, prime/radial mass, noncoprime correction, small primes, sieve error, level/Euler conversion and product BV & $E_{\rm Ssw}=o(1)$, once; product BV costs $O(L^{-5/2})$. \\
\end{longtable}
\endgroup

We take $H_L=I_{HL}$, the coordinate measure and $t_k$ of (5.30)--(5.32).
Set $W=W_{\rm norm}$ and put
\[
E_{\rm Ssw}(N)=\max\{0,S_{\rm sw}(N)/W-2I_{HL}\}.
\]
Then $-S_{\rm sw}/W\ge-2I_{HL}-E_{\rm Ssw}$, and
$I_{HL}\le H_{L,\rm upper}$ gives the sole main term $-2H_{L,\rm upper}$.
For fixed geometry and BV exponents, (5.43) gives
\[
\begin{split}
E_{\rm Ssw}\ll_{\varepsilon_A}{}&
L^{-1/3}+\frac{\log L}{L}+L^{-1}+L^{-2}
+\eta_P(N)+\eta_E(N)\\
&+L^{-5/2}+N^{-3/4}L^2+N^{-\alpha}L^3=o(1).
\end{split}
\]
Here $\eta_P,\eta_E\to0$ are the prime-box PNT/Euler errors of Section~5.5.
The radial and outer-box masses contribute $L^{-1},L^{-2}$.
By (5.40),
\[
R_{\rm nc}\le22X^+N^{-\alpha}\sum_{r<D}\frac{\mu^2(r)}{\varphi(r)}
\ll X^+N^{-\alpha}L.
\]
Since $X^+\ll N$ and $U_C\gg N/L^2$, its contribution is $O(N^{-\alpha}L^3)=o(1)$.
It enters $E_{\rm Ssw}$ once.
The $N^{-\alpha}L^3$ envelope covers (5.33) and (5.40), including source exceptions,
and $N^{-3/4}L^2$ covers the $36z$ endpoint.
The direct-BV conversions and the outer errors below are kept separate.

For Class II put $c_\sigma=\ell_0\sigma^2/4>0$ and
$C=C(\sigma,C_b,A',31,\mathcal C_{\rm src})$.
By \cite[Theorem~3.2, (3.13)]{Z1},
\[
|S_B|\ll_{\sigma,31,C_b,A',\mathcal C_{\rm src}}
XL^{-A'}+X^{1-c_\sigma}L^C.
\]
We sum the exterior coefficients $\le L^{K_0}$, use $X\ge N/L^{31}$ and divide by $W_{\rm norm}$:
\[
E_{\rm II}(N)\ll_\sigma L^{K_0+2-A'}
+N^{-c_\sigma}L^{K_0+C_{\rm pow}+2}.\tag{6.5}
\]
Here $C_{\rm pow}=C+31c_\sigma$ absorbs the fixed factor from
$X\ge N/L^{31}$.
Both terms are kept.
The cited proof contains GR, colouring, cross-copy coprimality, CRT,
zero mode, moments, Poisson, W-prime norm, Theorem T and endpoints
\cite[Theorem~3.2, \S\S3.5--3.9]{Z1}.
The actual-source H7 sum includes them once
\cite[\S3.10, (3.91)--(3.92)]{Z1}.
The outer errors are $O(N^{-\alpha}L^3)$ for rough exceptions,
$O(N^{-1/2}L^2)$ for squares/source factors dividing $N$,
and $o(1)$ for G4.
They enter (6.9) once, separately from O1.

\subsection{Selected endpoints and persistent costs}

We use the adverse endpoints of Appendix A.2--A.6. Terminating decimals are rational.
\begingroup\small
\setlength{\tabcolsep}{3pt}
\renewcommand{\arraystretch}{1.15}
\begin{longtable}{@{}p{0.22\textwidth}p{0.30\textwidth}p{0.40\textwidth}@{}}
Input & Certified value or enclosure & Use \\ \hline
\endfirsthead
Input & Certified value or enclosure & Use \\ \hline
\endhead
\multicolumn{3}{@{}p{\textwidth}@{}}{S7 at $\lambda=93/400$: canonical raw value $2.299260427561925683718\ldots$ in $[2.299209653294856,2.299402520476652]$; use the upper endpoint.} \\
\multicolumn{3}{@{}p{\textwidth}@{}}{S1 in D units: $3S_1/4\ge\underline D_1$, the lower endpoint of the S1 ball in Appendix A.2.} \\
\multicolumn{3}{@{}p{\textwidth}@{}}{Selected S2 endpoint: $d_{T8}>\underline d_2$.} \\
\multicolumn{3}{@{}p{\textwidth}@{}}{Independent S2 enclosure: $T8/4\in[1.62063469556829,1.62071322785677]$.} \\
\multicolumn{3}{@{}p{\textwidth}@{}}{S2 refinement: PF $0.000944946137955914$, corner $0.000058392441251611$, tail $0.00005$. The sum differs from the endpoint improvement by less than $2\cdot10^{-15}$.} \\
Switched integral & $H_{L,\rm upper}=0.0778506$ & $-2H_{L,\rm upper}=-0.1557012$. \\
Signed reference sum & $B_{\rm main}$ & $=1.83403742109459478063$. \\
Persistent errors & $F_{\rm pers}$ & $<0.000040767574227956903$; decomposition (6.6). \\
\end{longtable}
\endgroup

The two lower endpoints used above are
{\small
\[
\begin{aligned}
\underline D_1&=9.6765387489412441290823874931869409094969802156943270964
                 -6.00\cdot10^{-56},\\
\underline d_2&=1.61958135698908352559133404332105729220076714081650.
\end{aligned}
\]
}

The printed reference sum is the following thirteen-term sum of the
adverse endpoints in Appendix~A:
\[
\begin{aligned}
B_{\rm print}={}&\underline D_1+\underline d_2-\frac{U_{3,A}}2-\frac{U_4}4
-\frac{U_{5,\rm low}}4-\frac{U_7}4\\
&+\frac{L_{8,A}}4+\frac{L_{9,A}}4+\frac{L_{10}}4+L_{G1}
-\overline J_5-\overline J_6-2H_{L,\rm upper}.
\end{aligned}
\]
Here $\underline D_1$ is the lower endpoint of the S1 ball in Appendix~A.2,
and $\underline d_2$ is the selected S2 lower endpoint printed above.
The four $U$ endpoints and four $L$ endpoints are the adverse directed
rational endpoints printed in (A.7), (A.10), and Appendices~A.2--A.6.
The upper endpoints for the two remaining integrals are
{\small
\[
\begin{aligned}
\overline J_5&=\frac{0.03658166423843258896399704402764682632822246022922592990
                    +8.84\cdot10^{-57}}4,\\
\overline J_6&=1.3148581959157321533513247276782285310758770468886237619585.
\end{aligned}
\]
}
They come from (A.14) and (A.13), respectively.
The switched-integral endpoint is $H_{L,\rm upper}=0.0778506$.
Every summand is a printed rational endpoint; the sum uses no unrounded
analytic integral value.
Exact rational addition gives
$B_{\rm print}>1.83403742109641331345957914542194044932304169727869681073064$.
We take the smaller rational $B_{\rm main}=1.83403742109459478063$.
Thus $B_{\rm main}$ is obtained from the printed endpoint data.

We use the net-tail value $d_{T8}$ of (A.7).
Put $\delta_{T8}=1/20000$; this tail is already inside $d_{T8}$.
The persistent sum is
\[
\begin{split}
F_{\rm net}={}&0.000006484294227956901+0.00000575+0.00002853328\\
&+4.98\cdot10^{-54}+2\cdot10^{-24}+10^{-390}
<0.000040767574227956903.
\end{split}\tag{6.6}
\]
The order is the S6 increment, $g_{\rm BV}$, $g_{\rm AA}$,
the G1 reconciliation, $g_\tau$, $E_{1,fF}$.
By (1.27), $C_E\sigma W_{\rm AA}/4\le g_{\rm AA}$.
The exact BV decomposition is
\[
g_{\rm BV}=0.00000575
=\underbrace{0.00000004}_{E_{5,\kappa}}
+\underbrace{0.00000571}_{g_{R1}}.
\]
Here $E_{5,\kappa}=4\cdot10^{-8}$ is optional and conservative.
For comparison, Appendix~A.6 bounds the hypothetical S5-top increment
at $\kappa=10^{-6}$; the selected S5-top uses zero shift.
The selected S2 corner already uses $\kappa=10^{-6}$ and its increment
by (3.19) and (3.22).
The extra $g_{R1}\ge0$ is conservative. Both terms enter (6.6) once.

The four selector/cutoff/peel bounds sum to
\[
\begin{split}
\bigl(275009680+7.160+62500\\
\qquad+145664566.4567\bigr)\tau
&=420736753.6167\tau\\
&\le4.207367536167\times10^{-82}<2\times10^{-24}.
\end{split}
\]
The deliberately looser bound $4.201\times10^{12}\tau$ also covers this sum.
The allowance $2\times10^{-24}$ enters (6.6) once.
The independent S2 enclosure includes every tail:
\[
T8/4>d_{T8}.\tag{6.6a}
\]
Its lower bound improves $d_{T8}$ by about $0.001053338579$.
The input table gives the three improvements and the rounding error.
Sections~3.1--3.5 and~6.2 prove the actual-source arithmetic,
non-II, cutoff, AA37 and joint comparisons.

The equivalent pre-tail convention transfers only $\delta_{T8}$:
\[
\begin{aligned}
S2_{\rm pre}&=S2_{\rm net}+\delta_{T8},\\
B_{\rm pre}&=B_{\rm net}+\delta_{T8},\\
F_{\rm pre}&=F_{\rm net}+\delta_{T8}<0.000090767574227956903.
\end{aligned}\tag{6.6b}
\]
Thus $B_{\rm pre}-F_{\rm pre}=B_{\rm net}-F_{\rm net}$ exactly.
Here
\[
S2_{\rm pre}>1.61963135698908352559133404332105729220076714081650.
\]
The net value is smaller by $0.00005$,
and all other endpoints and error terms agree.
We use the net pair and subtract each error term once.
Then
\[
\begin{aligned}
B_{\rm main}-F_{\rm pers}&=B_{\rm net}-F_{\rm net}\\
&>1.83403742109459478063-0.000040767574227956903\\
&=1.833996653520366823727\\
&>1.83399665352036682372\\
&>1.78838848345960807361670991333>T.
\end{aligned}
\]
The strong outward margin exceeds $0.04560817006075875010$ D.
We retain the smaller safety margin
$D_{\rm lower}-T-g_{\rm AA}>0.04557963678075875010$.
This spends an extra $g_{\rm AA}=0.00002853328$ in the final vanishing-error
threshold; it is a reserve and is not charged a second time in (6.6).
The computation uses the full-precision (A.7), $H_{L,\rm upper}$ and the single G1 correction.
The displayed endpoint sum is $1.83403742109641\ldots$;
its excess over the conservative base is $1.82\cdot10^{-12}$.
Shorter intervals are outward summaries.
By (6.2) and (6.4b),
\[
\begin{gathered}
D_{\rm lower}=1.83399665352036682372,\qquad
D_{\rm src}(N)\ge D_{\rm lower}-E_{\rm van}(N),\\
E_{\rm van}(N)\longrightarrow0.
\end{gathered}\tag{6.3a}
\]

\subsection{Fixed parameters, then the limit}

First fix the source geometry and the classical constants.
Theorem~\ref{thm:sieve-mass} gives $W_{\rm AA}\le W_*=10^{14}$,
uniformly in tail-label counts.
The constant $C_E$ is absolute \cite[Theorem~2.6, (2.31)]{Z1}.
Choose
\[
0<\sigma\le10^{-10},\qquad
\sigma\le\frac{4g_{\rm AA}}{C_EW_*}\quad\text{if }C_E>0.
\]
Since $\eta_L/2500=4\cdot10^{-8}$ and AA37 has depth $R=1$
\cite[Theorem~2.6]{Z1}, the omitted restrictions
$\sigma\le\eta_L/2500$ and $\sigma<3^{-R}$ are automatic.
When $C_E=0$, only $0<\sigma\le10^{-10}$ is required.
Then $\tau=\sigma^9$ fixes the labels, factorial bounds and positive power savings.
Appendix A holds for $0<\sigma\le10^{-10}$ and $0\le\tau\le10^{-90}$.

By (4.20), with $A_0=53/256-\eta_L-\frac32\mu\tau$,
\[
\ell_c\ge\max(A_0-t_3,2t_3)\ge2A_0/3>1379/10000>\ell_0.
\]
Fix $A_{\rm raw}=3$, $C_b=334$ and then
\[
K_0=24C_b+26=8042,\quad A'=A_{\rm raw}+K_0+2=8047,\quad
A_{\rm work}=A'+4=8051.\tag{6.7}
\]
The G1 boundary exponent is $165-C_b/2$.
The condition $C_b>330$ gives decay, and the choice $334$ gives $O_\sigma(L^{-2})$.
The three slot-shell sums cost at most three logarithms, which leaves $A'$.
After counting the full Class-I family, we take $A_i\ge K_I+4$.
Choose the exceptional-cutoff and moment precisions by (4.67), including
$c_*\ge c_{\rm pre}+2$, so that all bounds in (1.29) hold.
The SW/BFI precisions are chosen next.
The width $334$ makes the source-wide boundary and freezing errors vanish.
The selected local precision gives the Class-II logarithmic term
$O_\sigma(L^{-3})$.
The resulting rates are
\begin{center}
\begin{tabular}{@{}lcc@{}}
Error & G1 & DIRECT \\\hline
Boundary & $O_\sigma(e^{163/L}L^{-2})$ & $O_\sigma(L^{-149})$ \\
Logarithm & $O_\sigma(L^{-168})$ & $O_\sigma(L^{-314})$
\end{tabular}
\end{center}
The upper G1 cut is $O_\sigma(L^{-331}+L^2/N)$;
Class I is $O_\sigma(L^{-2})$; the logarithmic term of (6.5) is $O_\sigma(L^{-3})$.
For fixed $v>0$ and finite $B$,
\[
\log(N^{-v}L^B)=-v\log N+B\log\log N\longrightarrow-\infty.
\]
Theorems~\ref{thm:repeated-primes} and~\ref{thm:sieve-mass}, with (6.5), give vanishing row errors.
The PNT/Mertens, BV/sieve and asymptotic AA37 errors are $o_\sigma(1)$.
The finite identities and the complete error table cover the sum.
Finally choose $N_0$. The choice includes all local theorem thresholds and requires $X\ge N/L^{31}$ and $\log X\ge L/2$. The endpoint slack conditions are
\[
\begin{aligned}
202(31\log L+1)&\le\sigma L,
&202\log(2^{24}/\varepsilon_A)&\le\sigma L,\\
202\cdot30\log2&\le\sigma L,
&\sigma L&\ge2^{130}(\log L+1).
\end{aligned}
\]
The last is the short-interval threshold (4.32a).
We also impose the GR roughness condition
$1/(\log\log X)^3\le\ell_0\sigma^2$ and all fixed short-interval moment and Brun--Bonferroni thresholds. The choice also includes (1.14) and the final slack in (1.23).
These are eventual conditions at fixed positive $\sigma$. For fixed $\kappa>0$, $C>0$, $d>0$ and real $B$, the quantity $CN^{-\kappa}L^B$ is at most $d$ when
\[
 L\ge\max\{1,(2\max(B,0)/\kappa)^2,
                         (2/\kappa)\log^+(C/d)\}.
\]
This follows from $\log L\le\sqrt L$. It shows directly that the
power errors in (1.30) are bounded by the last threshold choice.

For $N\ge N_0$ we also require
\[
E_{\rm van}(N)<\tfrac12(0.04557963678075875010).
\]
The weaker allowance is $D_{\rm lower}-T>0.04557963678075875010$.
The last threshold choice may be ineffective.

\subsection{From positive weight to a representation}

The results of Sections~3--5 hold for each fixed $0<\epsilon_A\le10^{-4}$.
Choose $\epsilon_A=10^{-4}$, $a=9/5$ and
\[
b=\frac{a-1}{a}-\epsilon_A=\frac{39991}{90000},\qquad
\frac{1-b}{b}=\frac{50009}{39991}.
\]
The exact threshold satisfies
\[
\begin{gathered}
1.78838848345960807361670991332<T=8\log(50009/39991),\\
T<1.78838848345960807361670991333,\\
D_{\rm lower}-T>0.04557963678075875010.
\end{gathered}
\]
We fix $\sigma$ before the limit.
By the definition of $\Sigma_W$ in (1.9), (6.1) gives
$\Sigma_W=W_{\rm norm}D_{\rm src}(N)$.
Theorem~\ref{thm:signed-assembly}, (6.2) and (6.3a) give
\[
\Sigma_W\ge
[4(1-\epsilon_A)D_{\rm lower}-o(1)]U_C.\tag{6.8}
\]
The pointwise sieve inequality subtracts the rough exceptional-set error.
Then the count is at least $[(1-\epsilon_A)D_{\rm lower}-o(1)]U_C$.
The fixed-$\epsilon_A$ G4 bound costs
$[8(1-\epsilon_A)\log((1-b)/b)+o(1)]U_C$.
Since $1-\epsilon_A>0$, $D_{\rm lower}>T$, and $U_C\to\infty$,
the integer count is eventually positive.

For a surviving semiprime $n=rq>\epsilon_A N$, with distinct primes
$r<q$ and $r<N^b$, one has $q>\epsilon_A N^{1-b}$ and
\[
(1-b)(a-1)-b=a\epsilon_A>0.
\]
Thus $r<q^{a-1}=q^{4/5}$ for
$N>\epsilon_A^{-(a-1)/(a\epsilon_A)}$.
For a prime survivor, $r=1$.
The positive integer count gives (1.1) in the quantitative form
\[
\operatorname{Good}(N)\ge
\big[(1-\epsilon_A)(D_{\rm lower}-T)-o(1)\big]U_C,
\qquad T=8\log(50009/39991).\tag{6.9}
\]
This proves Theorem~\ref{thm:goldbach}. $\square$
